\documentclass[11pt,a4paper]{article}
\usepackage[utf8]{inputenc}
\usepackage[italian]{babel}
\usepackage[margin=2cm]{geometry}
\usepackage{tabularx}
\usepackage{array}
\usepackage{graphicx}
\usepackage{xcolor}
\usepackage{booktabs}
\graphicspath{{images/}} 

\begin{document}
	
	\noindent
	\begin{minipage}{\textwidth}
		{\LARGE \textbf{Valutazione Euristica di Usabilità}} \\[0.5em]
		\textbf{Valutatore:} Mario Rossi \hfill \textbf{Data:} \today
	\end{minipage}
	\vspace{1.5em}
	\hrule
	\vspace{1.5em}
	
	\newcommand{\tabellaProblemi}[6]{
		\noindent
		\begin{tabularx}{\textwidth}{|X|}
			\hline
			\begin{minipage}[t]{\linewidth}\raggedright
				\vspace{4pt}
				\textbf{Problema N\textsuperscript{o}:} #1 \hfill \textbf{Grado di Severità:} #5\\[4pt]
				\textbf{Azione:} #2\\[6pt]
				\textbf{Domanda:} \par
				#3\\[6pt]
				\textbf{Descrizione:} \par
				#4
				\vspace{6pt}
			\end{minipage} \\ \hline
			
			\begin{minipage}[t]{\linewidth}\centering
				\vspace{4pt}
				\textbf{Screenshot:}
				\vspace{0.6em}
				
				#6
				\vspace{6pt}
			\end{minipage} \\ \hline
		\end{tabularx}
		\vspace{1em}
	}

	\tabellaProblemi{
		1 
	}{
		Il carrello non si aggiorna automaticamente dopo la modifica della quantità. 
	}{
		Il sistema fornisce un feedback visivo immediato sullo stato dell'applicazione? 
	}{
		Quando l'utente cambia la quantità di un articolo nel carrello da ``1'' a ``2'' e preme invio, il totale complessivo rimane invariato finché non si ricarica manualmente la pagina. Questo genera incertezza nell'utente circa la riuscita dell'operazione. 
	}{
		3 
	}{
		\textit{Nessuno screenshot allegato} 
	}
	
\end{document}